% from ega4-18-fr.tex
\oldpage[IV-4]{191}

  \item[c$'$)] Le morphisme $h$ est plat au voisinage du point $y$, et l'immersion $u$ est quasi-régulière au voisinage du point $y$.
\end{enumerate}
\end{proposition}

On a déjà prouvé dans (19.2.3) que $a)$ entraîne $b')$ sans hypothèse de finitude, et il est trivial que $b')$ entraîne $b)$. Montrons que $b)$ entraîne $c)$. On peut pour cela (la question étant locale sur $X$ et sur $S$) supposer que $Y$ est un sous-préschéma fermé de $X$ défini par un Idéal quasi-cohérent $\mathscr J$ de $\sh{O}_X$, et $u$ l'injection canonique. L'hypothèse que $\sh{O}_X/\mathscr J$ est $g$-plat entraîne que la suite
\[
  0\longrightarrow\mathscr J\otimes_{\sh{O}_s}\mathbf k(s)\longrightarrow\sh{O}_X\otimes_{\sh{O}_s}\mathbf k(s)\longrightarrow(\sh{O}_X/\mathscr J)\otimes_{\sh{O}_s}\mathbf k(s)\longrightarrow0
\]
est exacte ($0_{\mathrm I}$, 6.1.2), donc $Y_s$ est le sous-préschéma fermé de $X_s$ défini par l'Idéal $\mathscr J_s$ de $\sh{O}_{X_s}=\sh{O}_X\otimes_{\sh{O}_s}\mathbf k(s)$ qui s'identifie à $\mathscr J\otimes_{\sh{O}_s}\mathbf k(s)=\mathscr J/\mathfrak m_s\mathscr J$. Considérons une suite finie $(f_i)$ de sections de $\mathscr J$ au-dessus d'un voisinage de $y$ dans $X$, dont les images dans
\[
  \mathscr J_y/(\mathfrak m_s\mathscr J_y+\mathscr J_y^2)
\]
forment une base de ce $(\sh{O}_{X,y}/(\mathfrak m_s\sh{O}_{X,y}+\mathscr J_y))$-module (qui est libre en vertu de l'hypothèse $b)$). Puisque $X_s$ est localement noethérien, il résulte de (16.9.11), (16.9.5) et de l'hypothèse $b)$ que les images $(f_i)_s=f_i\otimes1$, sections de $\mathscr J_s=\mathscr J\otimes_{\sh{O}_s}\mathbf k(s)$ au-dessus d'un voisinage de $y$ dans $X_s$, forment une suite régulière engendrant la restriction de $\mathscr J_s$ à ce voisinage. Comme $\mathfrak m_s\sh{O}_{X,y}$ est contenu dans le radical de $\sh{O}_{X,y}$, et que dans les deux cas envisagés $\mathscr J_y$ est un idéal de type fini de $\sh{O}_{X,y}$, il résulte du lemme de Nakayama que les $(f_i)_y$ engendrent aussi $\mathscr J_y$; comme $\mathscr J$ est un Idéal de type fini de $\sh{O}_X$ dans les deux cas envisagés, il y a un voisinage $U$ de $y$ dans $X$ tel que les $f_i|U$ engendrent $\mathscr J|U$ ($0_{\mathrm I}$, 5.2.2). Le fait que $b)$ entraîne $c)$ est alors conséquence de ($0$, 15.1.16) lorsque $S$ et $X$ sont localement noethériens, et de (11.3.8) lorsque $g$ et $h$ sont localement de présentation finie; dans ce dernier cas, (11.3.8) prouve aussi l'équivalence de $c)$ et $c')$ (qui est triviale lorsque $X$ et $S$ sont localement noethériens ($0$, 15.1.11)). Enfin, si $c)$ est vérifiée, pour prouver $a)$, on peut encore se borner au cas où $Y$ est un sous-préschéma fermé de $X$ défini par $\mathscr J$, et par hypothèse, il y a une suite régulière $(f_i)$ de sections de $\mathscr J$ au-dessus d'un voisinage de $y$ dans $X$ telle que les $f_i$ engendrent $\mathscr J|U$ et que $h$ soit plat dans $U$; pour voir que $c)$ entraîne $a)$, il suffit encore d'appliquer ($0$, 15.1.6) lorsque $S$ et $X$ sont localement noethériens, et (11.3.8) lorsque $g$ et $h$ sont localement de présentation finie.

\begin{remarks}[19.2.5]
\label{IV.19.2.5}
\begin{enumerate}
  \item[(i)] Lorsque dans la condition $b)$, on supprime l'hypothèse que $h$ est plat au voisinage de $y$, la condition ainsi modifiée n'entraîne plus $a)$. Il suffit de prendre pour $S$ le spectre d'un anneau local qui n'est pas un corps, pour $s$ le point fermé de $S$, et $Y=X_s$.
  \item[(ii)] Supposons que $Y$ soit un sous-préschéma fermé de $X$ défini par un Idéal quasi-cohérent $\mathscr J$ de $\sh{O}_X$. On a prouvé au cours de la démonstration de (19.2.4) que, lorsque les conditions équivalentes de cette proposition sont remplies, alors, pour tout système $(f_i)$ de sections de $\mathscr J$ au-dessus d'un voisinage de $y$ dans $X$, tel que les images des $f_i$ dans $\mathscr J_y/(\mathfrak m_s\mathscr J_y+\mathscr J_y^2)$ forment une base de ce $(\sh{O}_{X,y}/(\mathfrak m_s\sh{O}_{X,y}+\mathscr J_y))$-module, il existe un voisinage ouvert $U$ de $y$ dans $X$ tel que les $f_i|U$ vérifient les conditions de la définition (19.2.1).
\end{enumerate}
\end{remarks}

